% !TEX program = xelatex
% lw-vs 中文主稿 —— QC-MDPC 首轮阈值时序研究
% 编译：xelatex main -> bibtex main -> xelatex main -> xelatex main
\documentclass[UTF8,a4paper,10pt,twocolumn]{ctexart}

\usepackage{amsmath,amssymb}
\usepackage{booktabs}
\usepackage{graphicx}
\usepackage{subcaption}
\usepackage{multirow}
\usepackage{array}
\usepackage{url}
\usepackage[hidelinks]{hyperref}
\usepackage{enumitem}
\usepackage{algorithm}
\usepackage{algpseudocode}

% ---- 版式 -----------------------------------------------------------------
\usepackage{geometry}
\geometry{a4paper,left=1.6cm,right=1.6cm,top=2.0cm,bottom=2.0cm,columnsep=0.65cm}
\setlength{\parindent}{2em}
\linespread{1.06}

% 图表标题：中文用「图 1」「表 1」，编号按节
\renewcommand{\figurename}{图}
\renewcommand{\tablename}{表}
\captionsetup{font=small,labelfont=bf,labelsep=quad}

% 跨栏浮动体
\newcommand{\dblfig}[3]{%
  \begin{figure*}[t]\centering
    \includegraphics[width=\textwidth]{#1}%
    \caption{#2}\label{#3}
  \end{figure*}}
\newcommand{\onefig}[4]{%
  \begin{figure}[t]\centering
    \includegraphics[width=\columnwidth]{#1}%
    \caption{#2}\label{#3}\vspace{#4}
  \end{figure}}

\newcommand{\cmark}{\ensuremath{\checkmark}}

\title{\vspace{-1.2em}\heiti\zihao{3} CLASSIC 类 QC-MDPC 比特翻转译码中首轮阈值时序\\对条件失败率的控制作用\\[0.35em]
\zihao{5}\mdseries 标定工作点下的密钥条件研究}

\author{\zihao{5} [作者姓名待补充]\\[0.2em]
\zihao{6} [单位、城市、邮编待补充]\\[0.2em]
\zihao{6} [通讯作者邮箱待补充]}

\date{}

\begin{document}
\maketitle
\thispagestyle{empty}

% ===========================================================================
\begin{abstract}
可靠译码是 QC-MDPC 类后量子密码的正确性前提，而其失败尾部同时受阈值调度与密钥结构支配。聚合解码失败率（DFR）无法回答一个更基本的问题：同一条译码实例上，阈值变化的\emph{时序配置}是否独立影响条件风险？我们定义 FRO-2——仅在首轮把经典阈值型比特翻转译码器（CLASSIC）的阈值提高 2，其后各轮不变——并在配对 $(H,e,s)$ 实例、以密钥为聚类单元的推断下评估。在 $(r,d,t)=(4801,45,92)$ 上，它在三个条件下把条件 DFR 降低 4.15--5.47 倍，全部密钥聚类区间排除 1；纯新密钥子集上最差组对数上界改善 0.687，通过预注册的 0.5 门槛。仪器化测量给出机制：错误首轮翻转下降约 $9.5\times$，而正确翻转仅下降 $2.7\times$，故即时残差增大，后续轮次在更干净的支持集上回收被推迟的修正。带轮次对照的消融表明，"偏移施加于何处"与"它是否持续"是两个彼此独立地起作用的性质。首轮是陡峭的双向窗口（加 2 改善 4.15 倍，减 1 恶化八倍）；而在同一批密钥上，仅在第 2 轮施加一次偏移使 DFR 降低 1.21 倍，同样的偏移若从第 2 轮起持续施加则摧毁译码器（0.017），尽管第 1 轮完全未被触及。效应无需重调即迁移至第二个工作点并在更大的码上复现，但对 $\kappa>d$ 的捕获倾向条件无保护作用，而一个近码字感知译码器恰在 FRO-2 有效处惰性、在 FRO-2 失效处起决定作用。在同一密钥池上把基准 DFR 再降低十四倍后，风险比为 3.71 [2.69, 5.57]——即保持稳定而非增强。在 BIKE Level-1 的码参数 $(12323,71)$ 上，把错误重量抬入可测带后首轮窗口同样是双向的（4.48 与 0.162）；在生产错误重量下测不到失败，但首轮决策边缘是可测的（$+2$ 把假翻转压低六倍）。最后，一个用首轮统计量预测其他偏移风险比的反事实模型在样本外失败——这些计数是处理后的——我们如实报告这一否定结果。
\end{abstract}

\vspace{1em}

% ===========================================================================
\section{引言}\label{sec:intro}

QC-MDPC 码以紧凑的稀疏校验结构支撑了 BIKE 这一基于编码的密钥封装机制，后者是 NIST 后量子标准化进程中的候选方案\cite{misoczki2013mdpc,bike2024spec,alagic2025nist}。其实用价值与一个困难的正确性问题相伴：迭代译码器结构极简，但残余失败概率位于无法直接测量的极端尾部。更关键的是，该尾部以安全分析所利用的方式依赖密钥：反应攻击、计时攻击与弱密钥分析均表明，译码行为会泄露或放大秘密校验矩阵的结构性质\cite{guo2016key,eaton2018timing,wang2023weak,gellersen2025}。这种依赖在\emph{可测}工作点上是否表现为密钥间失败率的离散，本身是一个经验问题；在本文主池上它没有被分辨出来——逐钥失败计数与同质泊松变异相容，我们把这一界如实报告而不是假定其不存在（第~\ref{sec:discussion} 节）。因此，真正的问题不只是"译码器平均多久失败一次"，而是"调度如何改变跨密钥与结构化错误条件下的条件风险"——无论残余离散最终是多少。

多数译码器研究优化全局纠错目标或平均 DFR。当密钥才是长期实验单元时，这类汇总并不充分：为同一把密钥生成的数千个错误样本，并不是数千个独立的密钥观测。一个调度可能改善聚合均值，同时把失败集中到少数密钥上；而若两次比较建立在不同的错误池上，则调度效应与样本难度将无法分离。由此得到三条设计要求：在同一条 $(H,e,s)$ 实例上配对两种配置、保留逐密钥结果、并以密钥为再抽样单元进行跨密钥推断。

本文考察的问题是：能否用一个\emph{刻意最小}的干预稳定译码器最早的决策。对经典阈值型比特翻转译码器（CLASSIC），我们定义 FRO-2 为"仅在第 1 轮把阈值加 2，此后使用未经修改的阈值"。该干预不含自适应状态、不含学习模型、不引入额外随机性。其理由是关于时序而非幅度的：初始综合征最重，保守的首轮决策可以在后续纠错动力学接管之前，阻止破坏性翻转。收益究竟真正绑定于首轮状态，还是仅仅取决于调度中某处存在偏移，正是下文带轮次对照的消融所要分离的问题。

本文的贡献有三。第一，建立了阈值干预的配对、密钥条件评估，以密钥而非单个错误抽样作为再抽样单元。第二，借助该协议我们表明：阈值偏移的安全性取决于其时序配置的两个彼此独立的性质——它触及哪一轮，以及它是一次性扰动还是持续位移；首轮是一个双向陡峭的窗口，而持续偏移后续轮次本身即是破坏性的。第三，把观测到的 DFR 下降与一个局部 UPC 尾模型联系起来——后者揭示出错误翻转被抑制的程度（约 $9.5\times$）显著强于正确翻转（约 $2.7\times$）——并刻画了该效应在工作点、阈值语义、迭代预算、密钥族以及一个捕获倾向的结构化错误边界上的迁移范围，同时把该干预定位在现代 QC-MDPC 译码器之间。第四，把测量推进到生产码参数：在 $(r,d)=(12323,71)$（即 Level-1 的原始分块）上，需要把错误重量提高约 $28\%$ 才能让基准失败率进入可测带；在该点上首轮窗口同样是双向的（风险比 4.48 与 0.162），而在生产错误重量本身测不到失败，可测的生产信号是首轮决策边缘与收敛代价——两者都在此量化，并附上规范自身译码器中首轮裕量的消融。此外我们报告一个否定结果：用首轮统计量预测其他偏移风险比的反事实模型无法跨偏移迁移，其原因我们也给出了诊断。

本文组织如下。第~\ref{sec:background} 节给出背景、相关工作与贡献边界。第~\ref{sec:model} 节建立首轮窗口的局部模型与三条可观测判据。第~\ref{sec:method} 节说明实验设计与统计协议。第~\ref{sec:results} 节报告结果。第~\ref{sec:discussion} 节讨论适用范围与局限。

% ===========================================================================
\section{背景与相关工作}\label{sec:background}

\subsection{QC-MDPC 比特翻转译码}

设 $H=[H_0\mid H_1]$ 为二元校验矩阵，由两个 $r\times r$ 循环块组成，块重为 $d$。对重量为 $t$ 的错误向量 $e$，译码从综合征 $s=He^{\top}$ 开始。并行比特翻转译码器为每个候选位置计算未满足校验数（UPC）$u_j$，翻转满足阈值规则的位置，随后更新综合征并重复，直至综合征归零或迭代预算耗尽。

本文全程使用的载体是\emph{经典阈值型比特翻转译码器}：即上文描述的并行比特翻转译码器，其阈值由当前综合征重量的二项模型导出\cite{misoczki2013mdpc,vasseur2021thesis}。按参考实现的命名，我们称其为 CLASSIC，其未经修改的调度记为 $C0$，并在全文作为配对对照。下文讨论的全部译码器——CLASSIC、BGF、BIKE-Flip、Backflip 与 BF-Max——都属于同一并行比特翻转家族，区别只在于阈值如何调度。译码器行为强烈依赖该调度：低阈值纠正更多真实错误位置，但也可能引入错误翻转；高阈值防止误翻转，却可能纠正不足而停滞。BGF 在仿射阈值上补充黑-灰两趟\cite{drucker2020shades}。BIKE-Flip 是 BIKE 规范 v5.2 中经过调优的生产译码器\cite{bike2024spec}，其裕量 $\delta$ 出现在\emph{每一轮}：以 $T=\max(f(w),F_i(w_0))$ 为阈值时，规范取 $F_1=f(w_0)+\delta$，$F_2,F_3$ 在 $f(w_0)$ 与下界 $(v{+}1)/2$ 之间过渡并同样加 $\delta$，$i\ge4$ 取 $(v{+}1)/2+\delta$；Level-1 上 $\delta=3$、固定七轮。已发表的机理依据是：失败源自第 2--3 轮，此时少数正确位置的计数器异常偏大，而由综合征导出的阈值已经下降\cite{sendrier2024seminar}。因此必须强调，FRO-2 \emph{不是}规范的调度：它是在二项阈值载体上对第 1 轮的一次性 $+2$，而规范的裕量是另一种调度形状中永久存在的一部分。Backflip 保存有限状态并允许基于存活时间的重翻转\cite{sendrier2020lowdfr}。BF-Max 则每轮选取计数器最大的位置并以随机方式打破平局\cite{baldelli2025bfmax}。近期的近码字感知方法通过更有针对性的状态与结构处理错误平台轨迹\cite{baldelli2026ncw,santini2026flip}。

FRO-2 并非一个新的译码器族，而是 CLASSIC 的单参数扰动，用于孤立首轮保守性的因果价值。正是这一狭窄范围，使得区分"时序阈值效应"与"一般性调度调参"成为可能。

\subsection{条件失败与弱结构}

QC-MDPC 译码器的失败概率以秘密密钥为条件。Guo 等证明译码错误会泄露距离谱信息\cite{guo2016key}，其后的计时与弱密钥研究进一步明确了攻击面\cite{eaton2018timing,wang2023weak,gellersen2025}。Nosouhi 等实现了 BIKE、量化了弱密钥对平均 DFR 的影响并提出密钥检查机制；该路线改变的是被接受密钥的总体\cite{nosouhi2023weakkey}。结构化密钥过滤可在密钥生成阶段降低暴露\cite{matthews2025combinatorial}，而译码器对策则试图让困难结构的影响变小。二者互补：过滤改变接受密钥的分布，调度改变密钥条件下的行为。

"弱密钥"一词涵盖互不相同的机制，不应作为单一类别使用。本文的 Type-I 构造强制某个校验块出现 12 个连续 1；在所测维度下这是对抗性罕见事件，用于对译码器施压，而非估计总体平均风险。浅交叠近码字条件注入的错误仅与选定近码字交叠 $\kappa=12$ 个位置，而 $\kappa<d=45$；由于到该近码字的距离仍在目标错误重量之上，该条件下捕获被构造性排除。因此它是\emph{浅交叠对照}，而非错误平台或 gathering 的代理。另有一组固定密钥批次的实验把 $\kappa$ 变化到 $d$ 以上，作为捕获倾向的结构化错误压力测试；其中观测到的失败本身并不构成进入动力学捕获盆的证据。

\subsection{贡献边界}

本文的定位需要明确说明，因为它限定了每一条结论的适用范围：

\begin{itemize}[leftmargin=1.4em,itemsep=0.15em,topsep=0.25em]
\item 本文是一项\emph{标定工作点下的译码器机制研究}，不是生产参数下的 DFR 定理，不是弱密钥恢复攻击，也不是 IND-CCA 安全性证明。
\item 主结果位于 $10^{-4}$--$10^{-3}$ 的可测压力点，距 BIKE 的密码学目标极远。本文不对失败率\emph{水平}作外推：生产参数一节报告的是在生产分块上、以抬高的错误重量测得的一个工作点，以及在生产错误重量本身的机制级统计量。
\item 与 BIKE-Flip、Backflip 的比较为语义与调参不对等的参照，不构成性能支配性主张。
\item 所谓的"捕获"条件一律表述为\emph{捕获倾向的结构化错误压力}，因为本文未记录轨迹终态距离或吸引域比例。
\item 阈值随机化在本设计空间内被否定，但这不证明所有随机化译码器无效——在位随机化与带状态方法差异显著\cite{baldi2021inplace}。
\item 首轮统计量不是经验证的预测器。形如 $P(\text{失败}\mid\text{首轮翻转计数})$ 的逐实例模型在基准臂内几乎完美地分类基准失败，却对每一个有利偏移都给出错误的符号；把偏移一并作为条件、并同时留出偏移，仍无法复现实测风险比。我们如实报告这一点，而不把首轮统计量呈现为预测装置。
\item 该生产参数点是生产的\emph{码}，不是生产的\emph{工作点}：只有错误重量被抬高，在其上测得的效应不构成关于设计重量处失败率的证据。
\end{itemize}

\subsection{相关建模工作}

有两条既有工作线与本文的机理研究重叠，并界定了本文的新意。第一条从轮级状态建模译码失败：Kim 与 Park 用（迭代序号，未满足校验数）的逻辑回归预测 LDPC 译码失败以提前终止\cite{kim2021logistic}；Paiva 与 Terada 以“首轮后剩余未纠正错误数”作为跨码长的小样本 DFR 代理\cite{paiva2022pickyfix}；Santini 等由“首轮未能减少错误数的概率”建立码特定的保守失败率模型\cite{santini2019conservative}；Annechini 等则由综合征重量分布估计 $(v,w)$-正则码的失败率\cite{annechini2024estimation}。第二条直接建模失败率本身：Sendrier 与 Vasseur 的（综合征重量×错误重量）马氏链\cite{sendrier2019dfr}、Tillich 的渐近律\cite{tillich2018dfp}、Vasseur 的两点外推框架\cite{vasseur2021thesis}、面向错误平台的密钥依赖马氏模型\cite{arpin2025errorfloor}，以及针对精确三轮并行译码器的闭式模型\cite{annechini2025closedform}——后者明确警告由代理译码器外推是危险的。上述工作均未把首轮统计量用于\emph{反事实}预测（即预测另一套阈值调度的风险比）；而我们自己的这一尝试在样本外失败，且原因可诊断（见第~\ref{sec:discussion} 节）。我们把该尝试及其失败如实报告，而不是把首轮统计量呈现为可预测装置。

% ===========================================================================
\section{首轮窗口的局部模型}\label{sec:model}

设 $U_{i,j}$ 为第 $i$ 轮中位置 $j$ 的 UPC 计数。在给定综合征重量与局部相关性的条件下，一阶模型为
\begin{equation}
\begin{aligned}
U_{i,j}\mid j\in\mathrm{err} &\sim \mathrm{Bin}(d,p_{1,i}),\\
U_{i,j}\mid j\notin\mathrm{err} &\sim \mathrm{Bin}(d,p_{0,i}),
\end{aligned}
\end{equation}
其中 $p_{1,i}>p_{0,i}$。对阈值 $\tau_i$，定义正确翻转与错误翻转的尾部
$q_{1,i}(\tau_i)$ 与 $q_{0,i}(\tau_i)$。FRO-2 仅把 $\tau_1$ 替换为 $\tau_1+2$，故两个尾部各自乘以
\begin{equation}
A_1=\frac{q_{1,1}(\tau_1+2)}{q_{1,1}(\tau_1)},\qquad
A_0=\frac{q_{0,1}(\tau_1+2)}{q_{0,1}(\tau_1)} .
\end{equation}
当阈值落在陡峭的错误翻转尾部、同时仍为真实错误保留质量时，模型预测 $A_0\ll A_1$：错误翻转被抑制的强度大于正确翻转。其预测轨迹是首轮残差更大，但进入级联或停滞盆的概率更低。

该近似以综合征重量为条件，并忽略 UPC 之间的相关性；它的用途是预测\emph{首轮的相对变化}，而非渐近 DFR。模型给出三条可观测判据：

\begin{enumerate}[label=(\roman*),leftmargin=1.8em,itemsep=0.12em,topsep=0.25em]
\item 错误首轮翻转的下降快于正确首轮翻转；
\item 持续施加的 $+2$ 偏移未必有益，即便它完全避开第 1 轮；
\item 在捕获倾向的 $\kappa>d$ 结构化错误条件下，小幅偏移未必有益。
\end{enumerate}

第~\ref{sec:results} 节的机制实验、含轮次对照的时序消融、以及高交叠压力实验分别检验这三条判据。

% ===========================================================================
\section{实验设计}\label{sec:method}

\subsection{证据分层与决策结构}

本文实验包含标定、候选筛选、留出确认、正式测量与预设敏感性分析五个阶段。图~\ref{fig:design} 概括了协议与证据结构；为便于读者判断每条结论的证据强度，表~\ref{tab:tiers} 明确列出四层证据的规模与结论权限。所有引用数字均可由冻结种子与配置从归档中复现，且第 1--2 层与第 3--4 层在正文中始终分别标注。

\begin{table}[t]
\centering
\caption{证据分层，按各层主导的规模列出；单个实验在正文中单独标注。第 3 层结果一律标注为探索性，不参与主结论。}
\label{tab:tiers}
\zihao{6}
\fontsize{6.2}{7.4}\selectfont
\setlength{\tabcolsep}{3pt}
\begin{tabular}{@{}clll@{}}
\toprule
层 & 角色 & 规模 & 结论权限 \\
\midrule
1 & 正式主实验及其对照 & 92--94 钥 $\times$ 3{,}000；留出 48 $\times$ 1{,}000 & 确证性 \\
2 & 预注册扩大确认 & 50 钥 $\times$ 1{,}000 & 确证性 \\
3 & 探索性消融 & 15--30 钥 $\times$ 1{,}000--2{,}000 & 仅探索 \\
4 & 原型工程对照 & 30 钥 $\times$ 500 & 可行性 \\
\bottomrule
\end{tabular}
\end{table}

整把密钥（而非错误实例）被分配至搜索、验证与测试各阶段。搜索阶段评估了 9 个调度：未修改载体、全局偏移、首轮偏移、两轮变体，以及掷币式阈值变体。逐次减半（successive halving）以最差条件组单侧 95\% DFR 上界为准则保留候选，并附加随机组非退化约束。最终选中的 FRO-2 仅把首轮阈值提高 2，其后没有任何阈值与 C0 不同。

\dblfig{../../figures/fig1_study_design.pdf}
{研究设计。(a) 配对译码协议：C0 与 FRO-2 处理同一条 $(H,e,s)$，真实错误仅用于生成实例与仪器化诊断；(b) 阈值时序：FRO-2 仅在第 1 轮将阈值提高 2，其后与 C0 完全相同；(c) 四层证据的规模与角色，第 3--4 层结果在全文中标注为探索性或原型级。}
{fig:design}

\subsection{工作点与条件组}

主可测工作点为 $P_1'=(r,d,t)=(4801,45,92)$，含两个循环块。其块长与块重取自 Misoczki 等为 80 位安全级推荐的参数集 $(n,r,\omega,t)=(9602,4801,90,84)$\cite{misoczki2013mdpc}，该集的每块块重为 $d=\omega/2=45$。我们保留 $(r,d)$，并把错误重量由推荐的 $84$ 提高到 $92$，使对照译码器的失败率高到可被直接仿真分辨。对 $t\in\{86,88,90,92,95\}$ 的标定扫描确认了这一取法：在 $t\le90$ 时 1{,}600 次标定抽样中未观测到 C0 失败，$t=92$ 时标定估计为 $3.1\times10^{-3}$。因此该点是一个为可测性而刻意过驱的正式推荐参数集，而非随手选取的代理点；它仍远离 BIKE 的生产参数，也不蕴含向 NIST 安全级的比例外推。

三个正式条件为：

\begin{itemize}[leftmargin=1.4em,itemsep=0.15em,topsep=0.25em]
\item \textbf{RAND}：独立生成的随机 QC-MDPC 密钥与均匀抽样的重量 92 错误；
\item \textbf{TYPE-I}：密钥被约束为在某一块内含有 12 个连续 1，配以均匀抽样错误；
\item \textbf{NCW-shallow}：密钥满足近码字判据，且错误被构造为与选定近码字交叠 12 个位置。由于 $12<d$，这是浅交叠对照而非捕获条件。
\end{itemize}

正式池含 92 把 RAND、94 把 TYPE-I 与 92 把 NCW-shallow 密钥，每钥 3{,}000 个错误，每臂共 834{,}000 个观测。其中纯新子集含 212 把未参与任何中期决策的密钥，每臂 636{,}000 个观测；该子集是预注册门槛的判定依据，全池结果作为敏感性分析报告。

迁移点为 $P_2=(3001,37,66)$，独立标定使 C0 再次具备可测失败。偏移在首轮严格保持 $+2$，未作重调。每条件 30 钥、每钥 3{,}000 错误，每臂每条件 90{,}000 个观测。

BIKE Level-1 使用原生参数集 $(r,d,t)=(12323,71,134)$，作为生产边界留出而非主要可测结果。生产尺度检查共两项：其一，偏移 $0$、$+2$、$+4$ 各在 RAND、TYPE-I 与近码字条件下运行一百万条实例，合计九百万次；其二，以官方 Level-1 常数实现的规范精确 BIKE-Flip 与 C0/FRO-2 在 16 钥、每条件 500 错误上比较。

预注册的扩大确认实验使用独立主种子 \texttt{20260912}，在 $P_2$ 与补充标定点 $P_3=(4001,41,82)$ 上各用 50 把新生成密钥、每钥 1{,}000 错误。六个调度共享 300{,}000 条唯一 $(H,e,s)$ 输入，产生 1{,}800{,}000 次译码臂运行。该实验与种子 \texttt{20260911} 的探索性矩阵完全分离，用于确认首轮位置效应与偏移幅度的条件依赖性。

\subsection{译码器语义与跨实现验证}

Python 实现以 $\mathbb{F}_2[x]/(x^r-1)$ 表示运算，同时包含参考与加速的 UPC 计算路径。单元测试比对 syndrome 构造、UPC 数值、翻转语义、生成器、失败分类与统计例程。BIKE-Flip 的仿射系数与整数阈值依据规范 v5.2 重建，并与已公布的 Level-1/3/5 调度对照验证。BF-Max 在计数器并列时使用伪随机打破平局，与已发表算法一致。

跨实现验证在两个层次上进行。其一，独立编写的动态距离 Python 路径观测到 11/20{,}000 次失败，C 参考为 136/200{,}000，两比例比较给出 $z=-0.68$，落在预设一致带内。其二，C 参考在每条件一百万条实例上复现了 FRO-2 在 RAND 与 Type-I 上的方向（分别为 1.59 倍与 2.01 倍），而 NCW 比较在该样本量下未分辨。两个实现由同一研究项目开发，故此检查检验的是实现一致性，而非独立实验室复现。动态阈值还使用了仿真特权信息（真实错误距离）。

主实验在阈值计算中使用静态目标错误重量 $t$，并以综合征归零判定成功。在一项独立的精确恢复审计中，我们追踪译码器所有翻转的奇偶性，把由此得到的估计错误与生成的错误比较；译码器本身不接收真实错误。审计覆盖全部 278 把正式 $P_1'$ 密钥、每臂 834{,}000 次译码，结果为 C0 精确恢复失败 795 次、FRO-2 失败 167 次，与零综合征失败总数完全一致，且两臂均未观察到零综合征错纠。这是被审计样本上的经验一致性，而非"零综合征"与"恢复原始错误"之间的逻辑等价。

算法~\ref{alg:paired} 给出配对译码与评估流程。C0 与 FRO-2 处理同一条 $(H,e,s)$，真实错误仅用于生成配对实例、标注诊断分析中的正确/错误翻转，以及在独立的精确恢复审计中比较估计错误。

\begin{algorithm}[t]
\caption{配对 C0/FRO-2 译码与评估}\label{alg:paired}
\zihao{6}
\begin{algorithmic}[1]
\Require 密钥 $H$，错误 $e$，阈值函数 $\tau(i,s)$，预算 $T$
\State $s_0 \gets H e^{\top} \bmod 2$ \Comment{实验驱动}
\For{$\mathrm{arm}\in\{\mathrm{C0},\mathrm{FRO\text{-}2}\}$}
  \State $s\gets s_0$
  \For{$i\gets 1,\dots,T$}
    \State $\delta \gets 2$ if $(\mathrm{arm}=\mathrm{FRO\text{-}2}$ and $i=1)$ else $0$
    \State $u \gets \mathrm{UPC}(H,s)$
    \State $F \gets \{j: u_j \ge \tau(i,s)+\delta\}$
    \State $s \gets s \oplus H\,\mathbf{1}_F^{\top} \bmod 2$
    \State 记录翻转位置与综合征重量
    \If{$s=0$} 记录成功并跳出 \EndIf
  \EndFor
  \If{$s\neq 0$} 记录失败 \EndIf
\EndFor
\State \Return 配对的 C0/FRO-2 结果
\end{algorithmic}
\end{algorithm}

对固定的 UPC 例程与迭代预算，FRO-2 与 C0 具有相同的渐近工作量与内存：稀疏 QC-MDPC 更新为 $O(Trd)$，译码器状态为 $O(r)$。其唯一的算法增量是第 1 轮的一次整数阈值偏移，不增加 UPC 遍历、不增加持久状态、不增加随机抽样。因此测得的时间差来自收敛轨迹的变化与实现差异，而非新的渐近操作。

\subsection{结果指标与统计}

主要结果为条件 DFR，即某条件某臂内失败译码数除以尝试数。由于 C0 与 FRO-2 处理同一条 $(H,e,s)$，个体结果保持配对。次要结果包括单侧 95\% Wilson 上界、逐密钥失败计数、收敛迭代数与失败分类。

主要跨密钥效应为风险比 $\mathrm{RR}=\widehat{\mathrm{DFR}}(\mathrm{C0})/\widehat{\mathrm{DFR}}(\mathrm{FRO\text{-}2})$。其 95\% 区间由以密钥为聚类、10{,}000 次重抽样、种子 \texttt{20260905} 的非参数聚类自助法给出；每次重抽样以有放回方式抽取密钥并保留该钥的全部错误。同时报告绝对 DFR 差、逐密钥配对失败计数的精确双侧符号检验，以及配对实例结果上的精确 McNemar 检验。符号检验与 McNemar 概率在三个正式条件上作 Holm 校正。另以 beta-二项敏感性模型估计调度优势比与钥内过度离散项；在这些数据的分辨率下，拟合的过度离散与零不可分辨（$p\approx0.49$），模型因此退化为合并风险比，给出的区间比密钥聚类自助更窄。故我们仅把它作为方向性敏感性检查报告，绝不用作独立的精度陈述。

本文在结果各节共报告了约五十个密钥聚类区间。Holm 校正只覆盖三个正式条件的符号检验与 McNemar 检验；风险比区间按构造是边际的——每个都是单一对照的 95\% 区间，而非该族的联合覆盖陈述，且未作校正。诸如"每个区间都排除 1"的表述应按单个对照理解，不应读作联合 95\% 主张；全文唯一作出多单元陈述之处是表~\ref{tab:temporal}，那里已明确说明。

对非零计数，DFR 以单侧 95\% Wilson 上界报告；对零失败，给出 rule-of-three 上界 $3/N$，且绝不把该估计解释为"风险为零"的证明。全文沿用统一的零计数约定，不引入隐含的连续性修正。

预设的操作门槛在最差条件组的单侧上界上以 $\log_{10}$ 尺度比较，要求至少 0.5 个 $\log_{10}$ 的改善且随机组不退化。这是决策规则，不是形式化的非劣性假设检验。原拟的异质性判据——逐密钥 DFR 的 P90--P10 跨度至少 3 个 $\log_{10}$ 单位——未被满足；观测跨度为 0.40--0.48 个 $\log_{10}$ 单位。

\subsection{实现与可复现性}

密钥、错误池、候选定义、阈值与随机选择均由主种子 \texttt{20260901} 经记录的 SHA-256 派生规则产生。逐密钥原始计数、配置文件、分析输出与 C 构建脚本均已归档。图源数据表由 \path{lw-vs/scripts/export_data.py} 直接从归档 JSON 导出，正文引用的每一个数字均可回溯至具体归档文件。

正式静态运行使用 Python 3.12、NumPy 2.4.6 与单进程；C 参考使用 Zig-CC 以 \texttt{-O2} 编译、未启用 AVX。归档未记录 CPU 型号与内存容量，故墙钟数值仅指示计算规模，不适合跨机器性能比较。

% ===========================================================================
\section{结果}\label{sec:results}

\subsection{时序结构：双向陡峭的首轮与持续偏移的代价}\label{sec:temporal}

阈值偏移的作用取决于其时序位置的两个彼此独立的性质：它触及哪一轮，以及它是一次性扰动还是对调度的持续位移。早先的设计无法把二者分开，因为全局偏移与首轮偏移的差别同时在于后续轮次和第 1 轮。为此我们补上了缺失的臂：从第 2 轮起持续施加、完全不触及第 1 轮的偏移，以及符号相反的首轮偏移。这两个臂相对于原始设计属于\emph{事后追加}——它们是在"全局对首轮"对比的混淆显现之后才确定的，并未预注册。我们按完整正式规模在三个条件上报告它们，但它们来自单一工作点；我们把这一点标为局限，而不把它们当作预注册确认来呈现。四种配置都在 $P_1'$ 的完整正式池上测量，C0 与 FRO-2 的计数与归档逐位一致，作为流水线自检（图~\ref{fig:temporal}d）。

首轮是一个陡峭的\emph{双向}窗口。把它的阈值提高 2 使 DFR 由 $1.022\times10^{-3}$ 降至 $2.464\times10^{-4}$（风险比 4.15 [3.29, 5.38]，92 把密钥中 79 把改善）。在同一轮降低 1 则在相反方向上同等剧烈：DFR 升至 $8.10\times10^{-3}$，风险比 0.126 [0.111, 0.142]，而且\emph{92 把密钥无一例外}都变差。保守的首轮决策有益、宽松的首轮决策有害，幅度相当且没有钥级例外——偏移的\emph{方向}与它的位置同等重要。

第二个性质与第一个彼此独立，而且在这个池上二者可以直接分开。从第 2 轮起施加 $+2$（完全不碰第 1 轮）把 DFR 抬升至 $6.08\times10^{-2}$，风险比 0.0168 [0.0149, 0.0187]，每一把密钥都变差。而\emph{仅}在第 2 轮施加 $+2$——同一批密钥、同一批错误池——反而把 DFR 降至 $8.48\times10^{-4}$，风险比 1.21 [1.07, 1.36]。两种配置的唯一差别是该偏移是否越过第 2 轮持续下去，而它们在三个条件上分处 1 的两侧：一次性时为 1.21--1.45，持续时为 0.012--0.021。由于两个臂共享同一密钥池，这一对比不带有"全局对首轮"比较所携带的工作点混淆。两个新臂在另外两个条件上同样复现：首轮负偏移的风险比为 TYPE-I 0.109 [0.099, 0.120]、NCW 0.101 [0.085, 0.120]；第 2 轮起持续偏移为 0.0210 [0.0192, 0.0229] 与 0.0118 [0.0100, 0.0137]，六个单元中没有任何一把密钥改善。因此具有破坏性的是对后续轮次调度的\emph{持续位移}，而不是"在第 1 轮之外存在偏移"这件事本身。

其余面板把这一结果放回更广的图景。在三个条件与两个工作点上，$+1$、$+2$ 与 $+3$ 的首轮偏移都把风险比抬升到 1.76--4.62，六个单元区间全部排除 1；而全局偏移把它压到 0.009--0.019（图~\ref{fig:temporal}a--c）。两种读法是一致的：在所测配置中，每一个一次性配置都改善译码器——首轮幅度大、第 2 轮幅度小——而任何持续配置无论从哪一轮开始都是破坏性的。

\begin{table*}[t]
\centering
\caption{时序消融：同一 $+2$ 偏移在不同轮次的作用（预注册扩大确认，50 钥 $\times$ 1{,}000 错误/钥）。风险比为 C0 DFR 除以该调度 DFR，$>1$ 表示优于 C0；粗体标出该行首轮偏移中的最大值，可见 FRO-2 并非每行最优。}
\label{tab:temporal}
\zihao{6}
\begin{tabular}{@{}llcccccc@{}}
\toprule
\multirow{2}{*}{工作点} & \multirow{2}{*}{条件} & \multirow{2}{*}{C0 DFR}
 & \multicolumn{3}{c}{首轮偏移} & \multirow{2}{*}{仅第 2 轮} & \multirow{2}{*}{全局 $+2$} \\
\cmidrule(lr){4-6}
 & & & FRO-1 & FRO-2 & FRO-3 & & \\
\midrule
$P_2$ & RAND        & $2.56\times10^{-3}$ & \textbf{2.56} & 2.37 & 1.91 & 1.13 & $0.013$ \\
$P_2$ & TYPE-I      & $3.14\times10^{-3}$ & 2.31 & \textbf{2.85} & 2.15 & 1.29 & $0.019$ \\
$P_2$ & NCW-shallow & $1.20\times10^{-3}$ & 1.76 & \textbf{3.00} & 2.50 & 1.36 & $0.009$ \\
$P_3$ & RAND        & $1.10\times10^{-3}$ & 2.29 & \textbf{2.62} & \textbf{2.62} & 1.34 & $0.010$ \\
$P_3$ & TYPE-I      & $1.70\times10^{-3}$ & \textbf{3.27} & \textbf{3.27} & 2.18 & 1.31 & $0.016$ \\
$P_3$ & NCW-shallow & $7.40\times10^{-4}$ & 3.70 & \textbf{4.62} & 3.08 & 1.54 & $0.009$ \\
\bottomrule
\end{tabular}
\end{table*}

这具有明确的机制含义。如果收益仅仅来自"把一个标定工作点移离阈值悬崖"，那么全局施加 $+2$ 应当改善译码器；它的大幅恶化排除了这一解释。全局偏移与首轮偏移同时在两个方面不同——它触及第 1 轮，并且位移了全部后续轮次——而上文的轮次对照表明，仅后者一项就足以摧毁译码器。表~
\ref{tab:temporal} 给出全部六个条件单元的点估计：被选中的 FRO-2 风险比为 2.37--4.62，全局偏移为 0.009--0.019；FRO-1 与 FRO-3 的跨度与之相近，说明关键在于\emph{首轮窗口本身}而非某一特定幅度。

值得指出，仅第 2 轮的偏移并非全无作用：四个单元（两个 TYPE-I 与两个 NCW）的区间排除 1，只是幅度仅约 1.1--1.5 倍，远小于首轮偏移。种子 \texttt{20260911} 的探索性矩阵（12 钥 $\times$ 400）只观测到 1.00--1.33 的信号，功率不足以分辨该效应，这也说明为何早期分析未能辨识它；此处以 50 钥 $\times$ 1{,}000 的预注册确认为准。

需要强调，上述均为六个边际区间而非未经调整的联合 95\% 覆盖陈述；最优偏移幅度也依赖于工作点与密钥条件——在 $P_3$/RAND 上 FRO-1 与 FRO-3 均略优于 FRO-2——因此 FRO-2 是一个\emph{经确认的代表性干预}，而非普适最优。

\dblfig{../../figures/fig4_temporal_dissociation.pdf}
{时序解离。(a)--(c) 三个条件上各调度相对 C0 的风险比，虚线为无效应基准 $\mathrm{RR}=1$；P2 与 P3 分别以实心圆与方形表示。首轮偏移（FRO-1/2/3）稳定位于 1 以上（1.76--4.62），仅第 2 轮偏移贴近基准但略高于 1（1.13--1.54），全局 $+2$ 塌缩至 $10^{-2}$ 量级。(d) $P_1'$ 完整正式池上的轮次对照，全部在同一批密钥与错误池上译码。首轮是双向的（$r_1$+2 改善 4.15 倍，首轮 $-$1 恶化 8 倍）；一次性与持续的区别在同一池内得到分离：仅第 2 轮施加 $+2$ 降低 DFR（风险比 1.21），而从第 2 轮起持续施加则摧毁译码器（0.0168），尽管第 1 轮完全未被触及。(a)--(c) 为 95\% 密钥聚类自助区间；(d) 的标注为相对 C0 的风险比。}
{fig:temporal}

\subsection{正式主实验}\label{sec:formal}

FRO-2 在三个正式条件下均降低失败数（表~\ref{tab:main}，图~\ref{fig:main}）。RAND 条件下失败数由 282/276{,}000 降至 68/276{,}000，对应 DFR $1.022\times10^{-3}$ 与 $2.464\times10^{-4}$（4.15 倍）；TYPE-I 由 372/282{,}000 降至 68/282{,}000，即 $1.319\times10^{-3}$ 降至 $2.411\times10^{-4}$（5.47 倍）；NCW 由 141/276{,}000 降至 31/276{,}000，即 $5.109\times10^{-4}$ 降至 $1.123\times10^{-4}$（4.55 倍）。

\begin{table}[t]
\centering
\caption{$(4801,45,92)$ 上的正式主结果。RR 为 C0 DFR 除以 FRO-2 DFR；区间为 95\% 密钥聚类自助区间；$p$ 值为 Holm 校正后的符号检验。}
\label{tab:main}
\fontsize{6.8}{8.2}\selectfont
\setlength{\tabcolsep}{2pt}
\begin{tabular}{@{}lcccc@{}}
\toprule
条件 & C0 失败/总数 & FRO-2 失败/总数 & RR [95\% CI] & Holm $p$ \\
\midrule
RAND & 282/276{,}000 & 68/276{,}000 & 4.15 [3.29, 5.38] & $4.9\times10^{-17}$ \\
TYPE-I & 372/282{,}000 & 68/282{,}000 & 5.47 [4.25, 7.28] & $3.0\times10^{-22}$ \\
NCW & 141/276{,}000 & 31/276{,}000 & 4.55 [3.21, 6.82] & $7.7\times10^{-12}$ \\
\bottomrule
\end{tabular}
\end{table}

所有聚类区间均排除 1，对应的绝对 DFR 差区间排除 0。在逐密钥计数上，FRO-2 在 233 把密钥上优于 C0、18 把劣于 C0、27 把持平（共 278 把）；配对实例上的精确 McNemar 检验经 Holm 校正后同样显著。在预注册的纯新密钥留出集上，最差组对数上界改善为 0.687，超过 0.5 的操作门槛，且 RAND 满足非退化约束。作为对照，正式全池的最差组改善为 0.679——两者方向一致，但\emph{纯新子集的 0.687 是门槛判定依据}，全池数值仅作敏感性报告，因为正式池是对中期池的扩展而非完全独立的新池。

纯新子集给出的风险比为 4.69 [3.58, 6.34]、5.52 [4.24, 7.52] 与 3.69 [2.53, 5.78]。beta-二项敏感性模型给出的 FRO-2/C0 优势比在全池为 0.241、0.183 与 0.220，区间均排除 1；在这些数据的分辨率下，拟合的额外钥级过度离散接近零。上述分析合起来表明，聚合结果既非由少数密钥造成，也非中期复用所致。

精确恢复审计进一步排除了"零综合征错纠"这一隐患：C0 的 795 次与 FRO-2 的 167 次精确恢复失败与零综合征失败总数逐一对应，两臂的零综合征错纠计数均为 0。

\dblfig{../../figures/fig2_main_results.pdf}
{正式主结果。(a) 三个条件上的条件 DFR 与单侧 95\% 上界；(b) 风险比 C0/FRO-2，横线为 95\% 密钥聚类自助区间，虚线为无效应基准；(c) 逐密钥失败数配对散点，虚线为对角线（无差异），点为密钥而非独立错误样本。}
{fig:main}

\subsection{机制：不对称抑制错误翻转}\label{sec:mechanism}

直接仪器化揭示了一个一致的首轮权衡（图~\ref{fig:mechanism}a）。C0 首轮平均翻转约 29 个位置，其中约 8 个是错误的；FRO-2 平均翻转约 9 个位置，其中约 0.8 个是错误的。因此\emph{错误}首轮翻转下降约 $9.5$ 倍，而\emph{正确}首轮翻转只下降约 $2.7$ 倍。由于 FRO-2 同时推迟了大量正确翻转，其首轮后的平均残差错误重量反而高出约 6 个位置。尽管如此，最终失败数在仪器化样本中由 RAND 的 21 降至 7、TYPE-I 的 25 降至 6、NCW 的 10 降至 1，平均迭代数增加约 0.25--0.28。

这一结果的含义是：FRO-2 并非在首轮解决了更多错误，而是\emph{刻意留下更多真实错误，同时注入远少的虚假错误}。后续轮次从更干净的支持集配置出发，回收了被推迟的修正。在被 FRO-2 挽救的子集中，基线平均产生约 16 次错误首轮翻转，而 FRO-2 为 1.4--2.2 次。

上述测量使用了真实错误支持集，因此是\emph{诊断性仪器}，并非部署译码器可获得的信息。这一区分很重要：该分析可以解释某个调度为何成功，但不暗示新的可实现决策信号。

\onefig{../../figures/fig3_mechanism.pdf}
{首轮机制。(a) 每轮译码的首轮翻转账目（RAND 条件，均值）：FRO-2 的错误翻转下降约 $9.5\times$，而正确翻转仅下降约 $2.7\times$，首轮后残差反而增加。计数为密钥与错误样本上的均值，单条译码实例仍翻转整数个位置。(b) 失败模式构成，按失败数归一化，分类采用预注册的 $\lceil T/4\rceil$ 窗口；30 个模式格子与归档的修正前分类器完全一致，原因是并行比特翻转在失败时必然耗尽预算。}
{fig:mechanism}{0.4em}

失败模式构成（图~\ref{fig:mechanism}b）以停滞为最大类别：C0 的失败模式计数在 RAND、TYPE-I 与 NCW 上分别为 198/282、257/372 与 101/141。归档计数由一个早期误翻转窗口为 $\lvert\mathrm{flips}\rvert/4$ 的分类器产生，而冻结定义为 $\lceil T/4\rceil$，因此我们在完整正式总体上按预注册定义重跑了分类。\textbf{30 个模式格子全部一致}，逐格对照表随发布包提供。我们可以给出原因，而不把它当成巧合：并行比特翻转在失败时从不提前终止，故每条失败轨迹都跑满预算，$\lvert\mathrm{flips}\rvert=T=20$，两个窗口在每条样本上都等于 5 轮——在这批数据上两个分类器是同一个函数。因此模式计数不再附带探索性保留，但它们是\emph{由构造等价而非新证据}得到确认的，我们如实这样陈述。直观测得的第一轮翻转账目仍是独立的诊断证据。

\subsection{迁移与敏感性}\label{sec:transfer}

在同一工作点 $(r,d,t)=(3001,37,66)$ 上报告了两项彼此独立的实验，不可混为一谈。本小节使用的是预注册的迁移运行：每条件 30 把密钥、每钥 3{,}000 个错误。表~\ref{tab:temporal} 使用的则是同一 $(r,d)$ 上的预注册扩大确认：50 把重新生成的密钥、每钥 1{,}000 个错误，并同时覆盖 $P_3$。两者在密钥、错误池与样本量上都不同，其 C0 的 DFR 也相应不同（RAND 上 $2.278\times10^{-3}$ 对 $2.56\times10^{-3}$）；二者互不包含。

冻结的 $+2$ 偏移无需重调即迁移至 $(3001,37,66)$（图~\ref{fig:scope}a）。RAND 的 DFR 由 $2.278\times10^{-3}$ 降至 $1.167\times10^{-3}$，TYPE-I 由 $3.000\times10^{-3}$ 降至 $9.333\times10^{-4}$，NCW 由 $1.200\times10^{-3}$ 降至 $4.778\times10^{-4}$；对应风险比为 1.95 [1.58, 2.42]、3.21 [2.57, 4.14] 与 2.51 [1.85, 3.46]。最差组对数上界改善为 0.384，小于主工作点但方向一致且统计上可分辨。

迭代预算敏感性把该效应与单纯的预算耗尽分离开（图~\ref{fig:scope}c）：在预算 7、10、20 下，RAND 与 TYPE-I 的风险比区间均排除 1，点估计为 1.72--3.17。预算 20 的区间较宽，是因为若干自助重抽样中 FRO-2 失败数为零；这一稀疏性如实报告而非被抑制。固定均匀误差的密钥族分析同样支持 FRO-2：随机、Type-I 与近码字可接受密钥批次的风险比为 4.17、4.44 与 4.33；而近码字可接受批次在均匀误差下的基线 DFR 仅为随机密钥的 1.04 倍，说明"仅凭准入条件"并不产生浅交叠对照效应。

在仿真特权的动态距离语义下，FRO-2 把 RAND 的 DFR 由 $7.38\times10^{-4}$ 降至 $4.63\times10^{-4}$、TYPE-I 由 $8.53\times10^{-4}$ 降至 $4.25\times10^{-4}$。1.59 倍与 2.01 倍的效果小于静态语义但分辨良好（$z=7.94$ 与 $11.98$）。NCW 由 $1.07\times10^{-4}$ 移至 $8.8\times10^{-5}$，统计上未分辨（$z=1.36$）。因此主结论并不依赖于静态载体较弱的假设，尽管效应幅度确实依赖语义。

\dblfig{../../figures/fig5_scope_and_boundary.pdf}
{迁移与边界。(a) 第二个标定点 $P_2$ 上的密钥聚类 95\% 区间，偏移未重调；(b) 交叠量 $\kappa$ 的捕获梯度，竖虚线为 $\kappa=d=45$。FRO-2（圆点）在 $\kappa\le d$ 时有益、一旦 $\kappa>d$ 即塌缩到无效应基准；而近码字感知译码器（方点）在 $\kappa\le d$ 时完全惰性、在 $\kappa=46$ 处转为决定性；(c) 迭代预算敏感性，RAND 与 TYPE-I 在所有所测预算下区间均排除 1；未生成 NCW 的预算消融数据，故无第三条曲线。}
{fig:scope}

该效应并不依附于单一工作点。在同一码上扫描目标重量时，FRO-2 在每一个能够分辨失败的工作点上都降低 DFR，且收益随译码器退化而单调缩小：$t=92$ 时风险比 7.00 [3.64, 23.5]、$t=94$ 时 2.82 [2.28, 3.60]、$t=96$ 时 1.87 [1.72, 2.04]；而在 $t\le90$ 处，早期的 30 钥扫描在 30{,}000 条中至多看到 1 次失败，无从分辨；低失败率一端现已在主池本身上分辨清楚（第~\ref{sec:lowdfr} 节）。扩大码长后，在 $(r,d,t)=(6007,56,100)$ —— 本研究最大的码，且 $t/r$ 低于任何标定点 —— FRO-2 把 DFR 由 $8.13\times10^{-3}$ 降至 $3.38\times10^{-3}$，风险比 2.41 [2.09, 2.82]。我们同时记录：我们同样报告完整扫描而非其中的一部分：共尝试了四个新的 $(r,d)$ 点、十二个（点，$t$）格子，其中六格落在可测带下方（30{,}000 次抽样中 C0 零失败）、六格落在上方。只有 C0 失败率在 13--54\% 的三格产生了可分辨的对比，且在那里收益随译码器退化而单调趋近于 1（同一码上 $t=98$、$102$、$106$ 的风险比为 1.77、1.08 与 1.01）。这些是刀刃状的工作点，可测窗口很窄；未命中本身就是结果的一部分。

\subsection{同一密钥池上的更低失败率点}\label{sec:lowdfr}

上一节的扫描把低失败率方向留作悬案：早期 $t$ 扫描在 $t\le90$ 处未能分辨出对比。因此我们在与主结果\emph{完全相同}的 92 钥密钥池与实例上重做 $t=90$，并采用预注册的停止规则（采样至 120 次基准失败或每钥 100{,}000 个错误，先到者止）；该规则在 1{,}840{,}000 条配对试验中的 130 次基准失败处触发，即每钥 20{,}000 个错误。基准 DFR 降至 $7.07\times10^{-5}$——为主标定点的十四分之一——风险比为 3.71 [2.69, 5.57]，60 把密钥改善、6 把变差、26 把持平（符号检验 $p=2.7\times10^{-12}$）。该区间与主标定点的 4.15 [3.29, 5.38] 大幅重叠，点估计相差约 $11\%$，因此在所测范围内，效应在失败率上保持\emph{稳定}而非随失败变稀而增强；区间宽度尚不足以排除中等幅度的变化，我们也不作此主张。跨池的 $t=92$ 扫描点（7.00 [3.64, 23.5]，另一个 30 钥池）的点估计落在主标定点区间之外，而两区间本身仍然重叠：池间点估计差异大于该范围内错误重量引起的差异——这与本文的密钥条件框架自洽。

\subsection{预规定偏移选择规则及其在留出工作点上的验证}\label{sec:rule}

前述结果会招致一个具体的质疑：$+2$ 究竟是对某个标定点的经验性调参，还是某个程序化过程的输出？为回答它，我们在跑任何实验\emph{之前}就冻结了一条选择规则，并只在一个此前从未用于本项目任何其他用途的工作点 $P_4=(3607,39,76)$ 上应用一次。该规则在 $P_4$ 上不运行任何译码实验，只依赖首轮 UPC 统计量。

设 $\tau_1$ 为 CLASSIC 首轮阈值，$q_c(\tau)$ 为首轮 UPC 值在 $\tau$ 处的尾部质量，分别对真实错误位置（$c=1$）与非错误位置（$c=0$）统计。把阈值提高 $\delta$ 每实例可避免 $(n-t)\,[q_0(\tau_1)-q_0(\tau_1+\delta)]$ 次错误翻转，代价是推迟 $t\,[q_1(\tau_1)-q_1(\tau_1+\delta)]$ 次正确翻转，故直接对交换率评分：
\begin{equation}
\begin{aligned}
\rho(\delta)&=\frac{(n-t)\,[q_0(\tau_1)-q_0(\tau_1+\delta)]}
                   {t\,[q_1(\tau_1)-q_1(\tau_1+\delta)]},\\
\delta^{*}&=\arg\max_{\delta\in\{1,2,3,4\}}\rho(\delta),
\end{aligned}
\end{equation}
并列时取较小偏移。该判据不含任何自由数值参数。一个更早的提法——按单位 $\delta$ 最大化不对称比 $A_1/A_0$——在正式运行前已被弃用，因为 $A_0\to0$ 会使其发散并总是选中最大偏移。

在 $P_4$ 上规则选出 $\delta^{*}=1$（图~\ref{fig:rule}a）。随后我们在 48 把新生成密钥、每钥 1{,}000 个错误上验证：C0、规则选出的 FRO-1 与冻结的 FRO-2 译码同一条 $(H,e,s)$ 实例（图~\ref{fig:rule}b）。C0 失败 40/48{,}000（DFR $8.33\times10^{-4}$）；FRO-1 与 FRO-2 各失败 8/48{,}000（DFR $1.67\times10^{-4}$），对应风险比 5.00 [2.64, 16.67] 与 5.00 [2.91, 12.00]，均排除 1，配对符号检验 $p=5.6\times10^{-6}$ 与 $3.0\times10^{-6}$。

这一结果的两层含义必须分开陈述。它\emph{支持}的是：首轮效应在一个未参与本研究任何开发环节的工作点上复现，且一个只依赖首轮 UPC 统计量的程序独立地得到了可用的偏移。它\emph{不}支持的是"该规则能在各偏移之间有效判别"：在 $P_4$ 上 $\rho$ 在 $\{1,2,3,4\}$ 上单调递减（0.628、0.486、0.409、0.373），故判据选中的是其范围内的最小偏移，且两个臂在该样本量下无法分辨。这条曲线说明的是\emph{最便宜的}偏移已经拿到了最好的可用交换率，而非 $+1$ 唯一最优。这与全文一贯的定位一致：FRO-2 是一个经确认的代表性干预，最优幅度依赖于工作点。

\onefig{../../figures/fig6_offset_rule.pdf}
{预规定偏移规则及其留出验证。(a) $P_4$ 上的交换率 $\rho$，由 12 把开发密钥的首轮 UPC 统计量算出，未运行任何译码器；规则选出 $\delta^{*}=1$。曲线单调递减，说明该判据识别的是"最便宜且足够"的偏移，而非在大偏移之间做判别。(b) 在 $P_4$ 上以 48 把新密钥 $\times$ 1{,}000 错误验证，三个臂译码同一条实例。柱为观测 DFR，竖线为单侧 95\% Wilson 上界；标注给出失败数/试次数与相对 C0 的密钥聚类自助风险比。}
{fig:rule}{0.4em}

\subsection{边界：捕获倾向的高交叠压力}\label{sec:boundary}

浅交叠 NCW 条件不应与本节的压力测试混同。我们固定一个近码字可接受密钥批次，令错误构造在均匀误差与 $\kappa\in\{12,46,60\}$ 之间变化（图~\ref{fig:scope}b）。FRO-2 在均匀与 $\kappa=12$ 条件下分别改善 3.54 [2.41, 6.00] 与 4.75 [2.00, 19.0] 倍。当 $\kappa=46>d$ 时，C0 在 85.31\% 的样本上失败、FRO-2 在 86.50\% 上失败，风险比 0.986 [0.981, 0.991]，轻微偏向 C0；当 $\kappa=60$ 时两者在全部 40{,}000 个样本上失败。

这构成一个清晰的实证范围边界：FRO-2 稳定的是可测可恢复条件下的早期阈值决策，而高交叠压力条件仍然高度易失败。面向错误平台或捕获恢复的方法处理的是另一种失败机制\cite{baldelli2026ncw,santini2026flip}。需要重申，这些观测表明在所测构造下失败率极高，但并不独立验证动力学捕获盆的存在。

一个结构不同的译码器表明，这条边界属于干预本身而非实验设定。我们实现了 Baldelli 等\cite{baldelli2026ncw} 的算法 2——识别出 syndrome 已到达近码字、并翻转该近码字全部支撑的比特翻转译码器——使用相同的 CLASSIC 阈值，并在同一条实例上运行。在 $P_1'$ 上，$\kappa<d$ 且捕获被构造性排除，该方法是\emph{完全惰性}的：其近码字步骤在全部 834{,}000 条实例中一次都未触发，失败数与 C0 逐位相同（282 / 372 / 141），平均迭代也完全相同（4.312、4.511、4.206）。而在捕获条件下对比十分鲜明（图~\ref{fig:scope}b）。在 $\kappa=46>d$ 处，C0 在 85.31\% 的实例上失败，而近码字感知译码器仅在 11.07\% 上失败，这一改善由触发率达 74.2\% 的步骤实现。在 $\kappa=60$ 处它也崩溃（成功率 0.76\%），说明即便是为此专门设计的结构性方法同样存在边界。在所有单元中，零综合征成功数与精确恢复数相等，故未观测到错纠——正是精确恢复审计使这一陈述可被检验而非被假定。

这一解离是对本文框架最清晰的支持。FRO-2 在早期阈值决策可恢复之处降低条件失败，且在 $\kappa=46$ 处近乎中性（风险比 0.986）；而为陷集设计的方法恰在 FRO-2 有效处惰性、在 FRO-2 失效处起决定作用。两种干预处理的是不同的失败机制，互不包含。

\subsection{负对照：阈值随机化}\label{sec:negative}

掷币驱动的阈值扰动未能改善条件风险。在一个边际固定阈值载体——BF30，即采用固定阈值 $b=30$ 的普通并行比特翻转，在筛选工作点 $P_1=(4801,45,84)$ 上——，对称抖动 $\{-1,0,+1\}$ 把 DFR 由 4.29\% 抬升至 19.67\%，仅正向抖动 $\{0,+1,+2\}$ 则在三个条件上抬升至 50.5--51.8\%：更低阈值造成过度翻转与发散，更高阈值造成翻转不足与停滞，安全带宽窄于扰动支撑集。

在 $P_1'$ 的 CLASSIC 上，同样的扰动方向一致，但该探针的样本量不足以分辨：RAND 为 2/2{,}400 对 7/2{,}400，TYPE-I 为 3/1{,}800 对 7/1{,}800，NCW 为 6/2{,}400 对 10/2{,}400，Fisher 精确检验下均不显著（Holm 校正后 $p\ge 0.54$）。因此我们把 CLASSIC 上的随机化结果报告为\emph{未分辨}而非证据，负对照的结论建立在固定载体上——那里的效应跨越一个数量级。在正式候选搜索中，两个掷币式变体均未通过筛选。这些结果在本设计空间内否定了随机化假设，但\emph{不}证明所有随机化译码器无效——在位随机化与带状态方法与本文的无结构阈值混合差异显著\cite{baldi2021inplace}。可得到的结论是：无结构的阈值混合不能替代定位良好的确定性调度。

\subsection{在现代译码器之间的定位}\label{sec:modern}

表~\ref{tab:modern} 把证据分为三层，而非形成单一排名：(i) 在完整正式池上对 C0、FRO-2 与规范公式 BIKE-Flip 的匹配静态语义比较；(ii) 使用官方 Level-1 常数的原生生产参照；(iii) C 实现动态距离规则下 Backflip 的语义参照。

\begin{table*}[t]
\centering
\caption{现代译码器比较与跨实现验证。DFR 为观测值；零失败单元给出 rule-of-three 上界。}
\label{tab:modern}
\zihao{6}
\begin{tabular}{@{}p{2.5cm}p{2.4cm}p{2.6cm}rp{1.5cm}p{4.4cm}@{}}
\toprule
比较层次 & 条件/参数 & 译码器 & 试次 & 失败数 & 解读 \\
\midrule
跨实现检查 & $P_1'$，动态语义 & Python/C CLASSIC & 20{,}000 / 200{,}000 & 11 / 136 & $z=-0.68$，一致 \\
\midrule
匹配静态语义 & RAND & C0 / FRO-2 / BIKE-Flip / BGF-C3 & 276{,}000/臂 & 282 / 68 / 83 / 287 & FRO-2 最低；BGF-C3 与 C0 未分辨\\
匹配静态语义 & TYPE-I & C0 / FRO-2 / BIKE-Flip / BGF-C3 & 282{,}000/臂 & 372 / 68 / 94 / 391 & FRO-2 最低；BGF-C3 与 C0 未分辨\\
匹配静态语义 & NCW & C0 / FRO-2 / BIKE-Flip / BGF-C3 & 276{,}000/臂 & 141 / 31 / 58 / 84 & BGF-C3 优于 C0 但低于 FRO-2\\
\midrule
原生 Level-1 & RAND/TYPE-I/NCW & 规范精确 BIKE-Flip & 8{,}000/单元 & 0 / 0 / 0 & $<3.75\times10^{-4}$；16 钥/单元 \\
原生 Level-1 & 同上 & C0 / FRO-2 & 8{,}000/单元 & 0 / 0 & 与 BIKE-Flip 未分辨 \\
\midrule
生产尺度留出 & 偏移 $0/+2/+4$ & C0 载体 & $10^6$/单元 & 全 0 & 上界 $<3\times10^{-6}$ \\
动态语义参照 & 三条件 & Backflip & $10^6$/单元 & 0 & 更强，但状态与迭代成本更高 \\
\bottomrule
\end{tabular}
\end{table*}

在官方 Level-1 常数下，规范精确 BIKE-Flip 同样在每条件 8{,}000 次试验中记录到零失败，并在约三轮内收敛。C0 平均迭代 2.05--2.12，FRO-2 为 2.40--2.56。BIKE-Flip 跨条件稳定的三轮行为体现了完全调优调度的价值。在无官方常数的标定点 $P_1'$ 上，由规范导出的系数给出 DFR $3.007\times10^{-4}$、$3.333\times10^{-4}$ 与 $2.101\times10^{-4}$：FRO-2 在 TYPE-I 与 NCW 上更低，在 RAND 上统计持平（McNemar $p=0.228$），但这是一次"被搜索过的偏移"对一个"域外导出的比较器"——不构成对 BIKE-Flip 的支配性证据。

在同一批实例与同一预算下，我们还运行了由参数自适应的 CLASSIC 阈值驱动的 GRAY\_BGF 结构（BGF-C3）。这是本研究的公平结构性参照：与规范精确变体不同，它使用的阈值调度与 $(r,d)=(4801,45)$ 相匹配。它在 RAND 与 TYPE-I 上与 C0 统计不可分辨（风险比 0.983 [0.836, 1.156] 与 0.951 [0.837, 1.086]），在 NCW 上优于 C0（1.679 [1.362, 2.111]）。FRO-2 在三个条件上均低于 BGF-C3，而 BGF-C3 的迭代成本高约 22\%（5.28、5.46、5.16 对 4.31、4.51、4.21）。我们也说明为何\emph{不}把规范精确的仿射变体用作该点的比较器：其公布的系数是 BIKE Level-1 的值，而项目自己的筛选记录写明 $d=45$ 上不存在调优常数。在此处施加该配置会在约简样本上产生 38--60\% 的失败率，这度量的是常数错配而非算法本身；我们报告这一边界，而不是把它当作基线。

在参照动态语义下，Backflip 在每条件一百万条实例上记录到零失败，给出低于 $3\times10^{-6}$ 的上界，在标定点上优于 C0 与 FRO-2，代价是额外的状态与更高的迭代成本。这一比较准确地定位了 FRO-2：它是一次极小的载体修改，而非所测最强的译码器。

\paragraph{生产边界检查。} 在原生 BIKE Level-1 上，官方 BIKE-Flip 与 C0/FRO-2 载体臂在原生比较中记录到零失败（每条件 8{,}000 条配对实例），九百万条实例的留出同样观测到零失败，只给出每单元 $<3\times10^{-6}$ 的 rule-of-three 上界；该留出是在 C 参考的动态距离语义下运行的，故其上界描述的是这一仿真特权载体而非静态载体，下文的收敛数据同样带有这一限定。可测的生产信号是\emph{收敛成本}：RAND 的两轮收敛率由 96.0\% 降至 68.4\% 再降至 2.2\%，TYPE-I 由 91.7\% 降至 60.2\% 再降至 1.5\%，NCW 由 93.6\% 降至 62.3\% 再降至 2.0\%（图~\ref{fig:modern}b）。

\paragraph{生产码参数。} 只抬高错误重量，就可以把分块参数本身带入可测机制。在 $(r,d)=(12323,71)$ 上，我们在 12 把密钥、每把 500 个错误上扫描基准失败率随错误重量的变化：失败率在 $t=170$ 处为 $1.67\times10^{-4}$、$t=172$ 处 $2.17\times10^{-3}$、$t=175$ 处 $1.5\times10^{-2}$、$t=180$ 处 $13.9\%$（更早一轮较粗的扫描补充 $t=195$ 处 99.9\%）。因此该分块参数下的可译码转变被限制在 $t$ 约 $15\%$ 的窗口内，而 $10^{-3}$ 水平位于 $t^{\ast}=172$，约高于生产重量 134 的 $28\%$。这是对“生产工作点距译码边缘有多远”的定量陈述，解释了为什么在该重量处测不到失败，而不仅是记录“没看到失败”。在 $t^{\ast}$ 上，生产分块参数下的首轮同样是双向的：32 把密钥、每把 1{,}500 条配对错误，$+2$ 偏移的风险比为 4.48 [2.97, 7.25]，单个 $-1$ 为 0.162 [0.119, 0.209]，基准在 48{,}000 条中失败 94 次。归档的生产运行只覆盖偏移 $0$、$+2$ 与 $+4$；负方向在此首次测量，且在能够分辨失败的重量处是决定性的。

在生产错误重量本身，所有臂都测不到失败（每臂 32{,}000 条配对试验，rule-of-three 上界 $9.4\times10^{-5}$）。在那里可测的是首轮决策边缘：基准平均翻转 101.2 个位置，其中 87.5 个是真错、13.8 个是假翻转；$+2$ 把假翻转压到 2.3（六倍），代价是缓置 25.5 个真错，而 $-1$ 把假翻转抬到 31.2。交换率——每缓置一个真错所避免的假翻转——在 $+2$ 处为 0.450，而标定点处为 0.546：同一结构，略更陡。可测的生产代价是收敛：$+2$ 平均需要 2.446 轮，基准为 2.054 轮，而 $-1$ 同样有代价（2.232 轮），故基准调度在三者中最快。规范自身的首轮裕量在其自己的译码器上表现相同：摘掉它使两轮收敛率由 0.09\% 升到 37.35\%，加倍则进一步压低该比例，而只在第 1 轮施加给出 6.14\%。四个消融臂各 8{,}000 次试验全部零失败，因此在 Level-1 的样本量下，可测的收敛代价与不可测的稳健性收益并存；我们如实报告两侧，不主张这笔交易划算。

\begin{table*}[t]
\centering
\caption{本研究新增的测量：抬入可测带的生产分块参数、生产错误重量本身、以及主池上的低失败率端。每个臂与其基准译码同一条 $(H,e,s)$ 实例；区间为按密钥聚类的 95\% bootstrap 区间。}
\label{tab:added}
\footnotesize
\begin{tabular}{@{}llrlrl@{}}
\toprule
设置 & 臂 & 基准 DFR & 臂失败/试验 & RR vs C0 [95\% CI] & 备注 \\
\midrule
$(4801,45,90)$，92 钥 $\times$ 20{,}000 & $+2$ & $7.07\times10^{-5}$ & 35 / 1{,}840{,}000 & 3.71 [2.69, 5.57] & 60/6/26 钥，$p=2.7\times10^{-12}$ \\
$(12323,71,172)$，32 钥 $\times$ 1{,}500 & $+2$ & $1.96\times10^{-3}$ & 21 / 48{,}000 & 4.48 [2.97, 7.25] & 生产分块参数 \\
$(12323,71,172)$，32 钥 $\times$ 1{,}500 & $-1$ & $1.96\times10^{-3}$ & 582 / 48{,}000 & 0.162 [0.119, 0.209] & 首次负方向测量 \\
$(12323,71,134)$，32 钥 $\times$ 1{,}000 & $+2$ & 0 / 32{,}000 & 0 / 32{,}000 & 未分辨 & 上界 $9.4\times10^{-5}$；仅有收敛代价 \\
$(12323,71,134)$，32 钥 $\times$ 1{,}000 & $-1$ & 0 / 32{,}000 & 0 / 32{,}000 & 未分辨 & 上界 $9.4\times10^{-5}$；仅有收敛代价 \\
\bottomrule
\end{tabular}
\end{table*}

\onefig{../../figures/fig7_modern_decoders.pdf}
{现代译码器比较与生产边界。(a) 标定点 $P_1'$ 上 C0、FRO-2 与规范公式 BIKE-Flip 的 DFR；RAND 条件上 FRO-2 与 BIKE-Flip 未分辨（McNemar $p=0.23$）。(b) 原生 BIKE Level-1 上两轮收敛率随首轮偏移的变化：每单元一百万次运行、零失败，因此该尺度上唯一可测的信号是收敛成本而非 DFR。}
{fig:modern}{0.4em}

\subsection{有界恢复的工程代价}\label{sec:engineering}

为把"极小的调度干预"与"带额外状态的译码器"分离开，我们运行了一组冻结的工程对照，包含 C0、FRO-2、STLR、SR-NCW 与 BGF-C3（表~\ref{tab:eng}）（图~\ref{fig:eng}）。STLR 在首个停滞轮后探测 UPC 最大的四个位置；SR-NCW 使用相同触发与预算，但在第三个条件（NCW-shallow）中把候选限制在严格优于规范的近码字支持集上。恢复用 UPC 向量在正常轮之后重新计算并计入成本。两个恢复臂均不接收真实错误支持集。

\begin{table*}[t]
\centering
\caption{两个标定点上的工程对照。所有报告单元均有非零失败；成本数值为 Python 参照测量，不构成硬件无关的速度主张。}
\label{tab:eng}
\zihao{6}
\begin{tabular}{@{}llrrrrr@{}}
\toprule
工作点 & 条件 & C0 & FRO-2 & STLR & SR-NCW & BGF-C3 \\
\midrule
$P_2$ & RAND        & 35/15{,}000 & 14/15{,}000 & 12/15{,}000 & 12/15{,}000 & 34/15{,}000 \\
$P_2$ & TYPE-I      & 36/15{,}000 & 24/15{,}000 & 23/15{,}000 & 23/15{,}000 & 64/15{,}000 \\
$P_2$ & NCW-shallow & 19/15{,}000 & 9/15{,}000  & 8/15{,}000  & 8/15{,}000  & 13/15{,}000 \\
$P_3$ & RAND        & 19/15{,}000 & 13/15{,}000 & 11/15{,}000 & 11/15{,}000 & 22/15{,}000 \\
$P_3$ & TYPE-I      & 25/15{,}000 & 9/15{,}000  & 9/15{,}000  & 9/15{,}000  & 24/15{,}000 \\
$P_3$ & NCW-shallow & 7/15{,}000  & 1/15{,}000  & 1/15{,}000  & 1/15{,}000  & 4/15{,}000 \\
\bottomrule
\end{tabular}
\end{table*}

STLR 与 SR-NCW 不只是接近。两者在全部六个单元上失败数相同（12/12、23/23、8/8、11/11、9/9、1/1），探测次数与触发率也相同；平均迭代仅在一个单元上相差 $7\times10^{-4}$（P2/NCW）。也就是说，近码字支持限制在全部 90{,}000 条配对实例上没有改变过任何一次决策，SR-NCW 相对 STLR 实质上是 no-op，而非一个独立的对照臂；我们如实这样报告。BGF-C3 每次译码使用 5.07--5.62 次 UPC 调用，而基于 CLASSIC 的各臂为 4.1--5.0 次，且并不一致优于 C0。这些区间有一个性质必须说明。当某个自助重采样抽到的密钥集合使第二个臂失败数为零时，其风险比为无穷；这些单元所用的分析代码把此类重采样\emph{丢弃}而非报告，在 $P_3$/NCW 单元中被丢弃的比例达 36.8\%。因此我们用归档的逐密钥计数重算了全部工程区间，保留无穷重采样，使受影响区间变为单侧。本段的读法不变——FRO-2/STLR 与 FRO-2/SR-NCW 仍在六个单元上全部包含 1——但仅按有限值给出的上端点低估了这些单元的不确定性，读者应视其为延伸至正无穷。

逐对区间显示：FRO-2/STLR 与 FRO-2/SR-NCW 的六个单元区间\emph{均未}分辨出额外收益（0/6），说明停滞触发的局部恢复与近码字支持限制在本文冻结设定下未带来可分辨增益；C0/BGF-C3 有 1/6 个单元排除 1，不存在一致的 DFR 优势。STLR 与 SR-NCW 每次译码额外产生约 0.002--0.02 次局部探测，触发率 0.05\%--0.51\%；其数据相关的触发与候选选择排除了常数时间主张。

工程结论是刻意受限的：首轮保守性是结构改动最小、最稳定的干预；停滞触发恢复在部分单元改变了 DFR/成本权衡；近码字支持限制未给出一致增益。BGF-C3 展示了在相同静态阈值下增加黑/灰子步的代价，但不是官方调优的 BGF 结果。

最后，我们在项目的规范 C 参考实现中测量了这两个载体：以 \texttt{zig cc -O2} 为 \texttt{x86\_64-windows-gnu} 编译、关闭 AVX、单线程、每次重复 50 万次译码，在机器空闲时取三次重复的最小墙钟。C0 每秒译码 7{,}134 次、FRO-2 为 6{,}563 次，吞吐比 1.087；平均迭代为 4.191 对 4.471（比 1.067），两者 P99 均为 6。因此编译实现在该干预上的代价约为 7--9\%，与 Python 侧的迭代计数一致，说明这一代价是真实的收敛效应而非解释器开销。两条限定约束该测量。其一，该参考使用动态距离阈值语义，阈值可参考到真实错误的当前距离；这是仿真特权，并非主实验的静态语义，故只解释 C0 与 FRO-2 的比值，不解释绝对 DFR。其二，同一二进制中的 GRAY\_BGF 构建使用上文所述的错配仿射常数，其 13.2 的平均迭代与 47.9\% 的预算耗尽率度量的是该错配而非 BGF 的代价；我们给出该数字只是为了让这一点明确。

\onefig{../../figures/fig8_engineering_tradeoff.pdf}
{工程成本--DFR 权衡。(a) 每个点为一种方法在一个工作点与条件单元上的观测：横轴为每次译码的 UPC 通过次数，纵轴为 DFR；(b) 各方法对在六个单元上的风险比均值与 95\% 区间，右侧标注区间排除 1 的单元数。FRO-2 相对 STLR 与 SR-NCW 的六个区间均未分辨出额外收益（0/6）。}
{fig:eng}{0.4em}

% ===========================================================================
\section{讨论}\label{sec:discussion}

\subsection{证据支持什么}

在主要正式池、纯新子集、第二个标定点、两种译码器语义、多个迭代预算与三个密钥族上，首轮 $+2$ 偏移总体降低所测 CLASSIC 载体的条件失败率。主实验与预注册 $P_2$ 迁移实验中的所有密钥聚类区间都超过 C0/FRO-2 风险比 1，且多数密钥单独改善；动态语义的 NCW 条件与若干稀疏工程单元仍统计未分辨。

更重要的是这一效应如何分解。全局偏移之所以有害，并不是因为它触及第 1 轮——一个只从第 2 轮起施加偏移、完全不碰第 1 轮的臂同样具有破坏性——而是因为它位移了每一个后续轮次，使被推迟的修正无法被回收。对称地，首轮是一个双向窗口：宽松的首轮决策与保守的首轮决策方向相反、量级相当。因此机制命题是：一次\emph{单次}、位置恰当的早期决策改变了进入后续迭代时的状态，而后续调度无法承受持续位移。

机制测量细化了这一解释。FRO-2 在严格意义上更保守：它做出更少的首轮翻转、远少的错误翻转，并留下更多真实错误。译码器因此可以在"即时残差重量看起来更差"的同时改善最终成功率。这一观测提醒：不应把早期错误削减本身当作调度目标；合适的目标必须同时考虑残差支持集的\emph{质量}与后续轮次的回收能力。

对首轮统计量的反事实使用失败了，而原因本身重要。由于首轮翻转计数与首轮后综合征重量可以从初始状态与候选偏移算出、无需运行该偏移的译码器，自然会想到在基准臂上拟合 $P(\text{失败}\mid\text{首轮统计量})$ 并用它预测其他偏移。在主池上，这样的模型在其自身臂内几乎完美地分类基准失败（留出密钥 AUC 0.995），却对每一个有利偏移都给出错误的符号——留出密钥上的预测风险比为 0.79、0.16 与 0.04，而实测为 3.46、3.86 与 3.97（$+1$、$+2$、$+3$）。把偏移一并作为条件、并同时留出密钥与偏移，仍然失败（2.18、1.86 与 0.83，对同样的实测值，且最大偏移的方向预测错误）；把同一个模型用到另外三个归档配置上——同一码的 $t=94$ 与 $t=96$、以及更大的 $(6007,56,100)$ 码——预测值为 0.21、0.27 与 0.33，而实测为 2.82、1.87 与 2.41，无一落在实测区间内。诊断是\emph{处理后变量}的条件化问题：较少的真错翻转在基准下意味着“实例本身难”，在抬高的阈值下却意味着“被干预推迟”，因此条件失败概率无法跨偏移迁移。这是 \cite{annechini2025closedform} 所警告的危险的一个实证实例，也界定了首轮统计量能支持什么：聚合交换率解释了保守首轮为何可能有益、并在所有所测点上与符号一致，但没有得到经过验证的定量预测器，我们也不作此主张。

有一个反驳需要正面回答。生产译码器在全部轮次都保留裕量 $\delta$，而本文却报告持续偏移会摧毁译码器。两句话说的都是对\emph{特定}调度的扰动。本文的持续偏移臂位移的是一个由当前综合征重量经二项式模型导出阈值的调度的后续轮次：该载体的每一轮都各自经过标定，同时位移全部轮次会破坏这一标定。规范的调度形状不同——其后段迭代以下界 $(v{+}1)/2$ 为底，裕量从第 1 轮起就按构造存在，其标定对象是整个调度而非每一轮单独标定。因此正确的读法是：对已标定的逐轮调度做\emph{中途}位移是有害的，而不是“任何含常数裕量的调度必然失败”。

负对照同样具有信息量。随机阈值选择改变了轨迹的条件混合，但并不自动降低密钥依赖或最差组风险；当可行阈值带很窄时，混合调度只是在采样失败悬崖。相比之下，确定性时序利用的是初始与后续综合征状态之间可重复的不对称性。

\subsection{适用范围与安全相关性}

本文实验测量的是译码器层面的条件风险，不展示实用密钥恢复攻击或 IND-CCA 安全性证明，也不否定 BIKE 的隐式拒绝保护。DFR 变化只有在部署参数、拒绝语义、常数时间约束与错误反馈接口都与实验匹配时，才可能转化为系统层面的可靠性或安全后果。FRO-2 仅改变首轮阈值计算且无需持久状态，因而是对 CLASSIC 类译码器的最小算法修改；但它仍带来可测的迭代与早期收敛代价。在可用的调优 BIKE-Flip 或 Backflip 面前，本文没有证据支持用 FRO-2 替换它们。

外推边界同样重要。主结果的 DFR 位于 $10^{-4}$--$10^{-3}$ 的标定压力点，远高于 BIKE 的密码学目标，全文不作生产尾部外推。所测高交叠结构化错误条件对 C0 与 FRO-2 都具有极高失败率，这些结果不构成一般性捕获盆的证据，该干预也未被证明是错误平台防御。生产尺度的零失败运行只排除了高于已实现分辨率的退化，不能证明收益。RAND 主要支持所抽样的随机密钥生成器；TYPE-I 是对抗性结构探测；NCW-shallow 是同时改变密钥准入与错误构造的浅交叠对照。固定均匀误差的近码字族只提供与 gathering 相邻的结构性证据，而非精确的 gathering 密钥结果。

\subsection{局限}

主要结论依赖两个为可测失败而选择的工作点，第三个标定点仅用于补充敏感性与扩大确认。其余局限包括：候选搜索空间由研究者定义；补充单元仅有 15--30 把密钥；跨实现检查在同一研究项目内完成，而非由独立实验室执行；首轮标注使用真实错误支持集，属诊断而非可部署信息。

三条边界陈述专属于本轮新增结果。其一，生产参数点是生产\emph{码}在抬高错误重量下的结果，而非生产工作点；该工作点由括号二分选出，其样本是主实验实例的子集；在其上测得的效应不构成关于设计重量处失败率的证据。其二，生产错误重量处的统计量是不解码即可计算（对反事实偏移而言）的机制测量，且该重量下的收敛比例差异是我们仅有的、统计上分辨清楚的生产信号。其三，在主池上，逐钥失败计数与同质泊松变异相容（方差/均值 0.98；密钥间失败率变异系数的 95\% 上界约 30\%），且两臂的逐钥失败集合几乎不相关，故本文的“密钥条件”描述的是\emph{推断方式}，而不是已证实的密钥级失败率离散。

未来工作应包括：增加独立标定的工作点或半解析尾部界；实现精确定义的 gathering 密钥构造；获得外部盲复现；在匹配的语义与预算下报告运行时间、UPC 计数、时序分布、常数时间约束与内存访问行为；以及为局部恢复原型给出显式的结构分数并验证其触发是否确实对应近码字结构。

% ===========================================================================
\section{结论}\label{sec:conclusion}

把首轮阈值提高 2、其后保持 CLASSIC 调度不变，在主标定点上把测得条件 DFR 降低约 4.1--5.5 倍。该方向在密钥聚类推断、纯新密钥分析、第二个标定点与参照语义下均持续存在。直接测量解释了其中的权衡：FRO-2 抑制错误首轮翻转，接受更大的即时残差，并依赖后续轮次完成回收——错误翻转下降约 9.5 倍，而正确翻转仅下降约 2.7 倍。

同一批实验也界定了该干预\emph{不}做什么。阈值随机化在所测机制下无效或有害；高交叠结构化错误条件仍然高度易失败，而近码字感知译码器才是那里的合适工具；在生产错误重量上，可用样本量对任何所测偏移都未能分辨出失败，因此未建立任何收益——在那里建立起来的是首轮决策边缘与收敛代价，而用首轮统计量构建的偏移效应反事实模型在样本外失败；调优的 BIKE-Flip 与带状态的 Backflip 仍是更合适的性能参照。

所支持的结论不是普适的生产优势，而是一条包含四部分的局部原理：在可测的 CLASSIC 工作点上，首轮的一次性保守偏移可以稳定部分译码轨迹；首轮窗口是双向的，因此宽松的首轮决策同样举足轻重；而持续偏移后续轮次会摧毁译码器，无论第 1 轮是否被触及；并且该原理在生产分块参数上依然成立——一旦把错误重量抬入可测带，那里的窗口同样是双向的，而生产重量本身的失败率仍不可及。所测高交叠结构化错误条件仍然易于失败。密钥相关的译码器研究应做到在相同实例上配对结果、在密钥层面切分、把密钥当作聚类、报告零失败的上界，并对边界条件做显式检验。

% ===========================================================================
\section*{数据与代码可用性}
\addcontentsline{toc}{section}{数据与代码可用性}
实现、冻结配置、种子派生规则、逐密钥原始计数与分析代码归档于 \url{https://doi.org/10.5281/zenodo.22296719}，开发镜像见 \url{https://github.com/Leekj12410/qc-mdpc-experiment}。归档覆盖正式主实验与轮次对照运行、迁移与扩大确认、机制、捕获与迭代消融实验、精确恢复审计，以及 BIKE-Flip、BGF 与近码字感知译码器的对照。本稿自身的产物随同一发布提供：图源数据表、从归档计数重新生成这些表的脚本，以及把正文引用的每个数值与这些表逐项比对的校验脚本。本文的每个数字都可沿该链条追溯到归档文件；校验脚本一并提供，使这一映射可被重新核验，而非只能被信任。

% ===========================================================================
\section*{作者贡献}
\addcontentsline{toc}{section}{作者贡献}
[待补充：作者姓名、单位与 CRediT 角色分工。]

\section*{利益冲突声明}
\addcontentsline{toc}{section}{利益冲突声明}
[待补充：作者利益冲突声明。]

% ===========================================================================
\bibliographystyle{IEEEtran}
\bibliography{refs}

\end{document}
